\documentclass[11pt]{article}

\usepackage[T1]{fontenc}
\usepackage{lmodern}
\usepackage{microtype}
\usepackage[a4paper,margin=25mm,headheight=14pt]{geometry}
\usepackage{xcolor}
\usepackage{enumitem}
\usepackage{booktabs}
\usepackage{tabularx}
\usepackage{array}
\usepackage{fancyhdr}
\usepackage{titlesec}
\usepackage{hyperref}
\usepackage{lastpage}
\usepackage{xspace}

\definecolor{navy}{HTML}{17324D}
\definecolor{teal}{HTML}{1E6F78}
\definecolor{lightblue}{HTML}{EAF2F7}
\definecolor{lightgray}{HTML}{F3F5F7}
\definecolor{rulegray}{HTML}{C8D1D8}
\definecolor{textgray}{HTML}{3F4A54}

\hypersetup{
  colorlinks=true,
  linkcolor=navy,
  urlcolor=teal,
  pdftitle={Recommendation Regarding arXiv Dissemination of the ALCI_RCC5 Overview},
  pdfauthor={GPT-5.6 Pro, OpenAI},
  pdfsubject={Publication recommendation},
  pdfkeywords={ALCI RCC5, description logic, arXiv, publication recommendation}
}

\pagestyle{fancy}
\fancyhf{}
\fancyhead[L]{\small\color{textgray}arXiv Publication Recommendation}
\fancyhead[R]{\small\color{textgray}$\mathsf{ALCI}_{\mathsf{RCC5}}$ overview}
\fancyfoot[L]{\small\color{textgray}20 July 2026}
\fancyfoot[R]{\small\color{textgray}Page \thepage\ of \pageref{LastPage}}
\renewcommand{\headrulewidth}{0.4pt}
\renewcommand{\footrulewidth}{0pt}

\titleformat{\section}
  {\large\bfseries\color{navy}}
  {\thesection.}{0.55em}{}
\titlespacing*{\section}{0pt}{1.4em}{0.55em}

\titleformat{\subsection}
  {\normalsize\bfseries\color{teal}}
  {\thesubsection}{0.5em}{}
\titlespacing*{\subsection}{0pt}{1.0em}{0.35em}

\setlist[itemize]{leftmargin=1.45em,itemsep=0.28em,topsep=0.35em}
\setlist[enumerate]{leftmargin=1.7em,itemsep=0.32em,topsep=0.35em}
\setlength{\parindent}{0pt}
\setlength{\parskip}{0.58em}
\renewcommand{\arraystretch}{1.18}

\newcommand{\logic}{\ensuremath{\mathsf{ALCI}_{\mathsf{RCC5}}}\xspace}
\newcommand{\fboxtext}[1]{%
  \begin{center}
  \setlength{\fboxsep}{10pt}%
  \colorbox{lightblue}{\parbox{0.91\textwidth}{#1}}
  \end{center}
}

\begin{document}

\begin{center}
  {\color{navy}\rule{\textwidth}{1.1pt}}\\[1.1em]
  {\LARGE\bfseries\color{navy}Recommendation Regarding arXiv Dissemination}\\[0.55em]
  {\large The final overview manuscript on satisfiability in \logic}\\[1.05em]
  \begin{tabular}{@{}rl@{}}
    \textbf{Manuscript assessed:} & \emph{Can Agentic AI Settle a 20-Year-Old Description Logic Problem?}\\
    \textbf{Author:} & Michael Wessel\\
    \textbf{Assessment date:} & 20 July 2026\\
    \textbf{Prepared by:} & GPT-5.6 Pro (OpenAI)\\
  \end{tabular}\\[1.0em]
  {\color{navy}\rule{\textwidth}{1.1pt}}
\end{center}

\fboxtext{%
\textbf{Recommendation: post the manuscript to arXiv.} It warrants public dissemination as an \textbf{overview, technical status report, and durable research handoff}. It should not be framed as a claimed solution of the open decidability problem, and the final revision already makes that distinction sufficiently clear.}

\section{Scope of this recommendation}

This assessment applies the standards appropriate to an overview paper. The manuscript is not required to reproduce all underlying proofs in self-contained form. Its principal obligations are instead to report the state of the problem accurately, distinguish certified results from theorem-level and empirical claims, provide usable pointers to the underlying work packages, and preserve enough of the research history that another group can resume the investigation without reconstructing it from scratch.

The recommendation concerns \emph{arXiv dissemination}. It is not an assertion that the open problem has been solved, nor a substitute for specialist peer review of every theorem, formal artifact, implementation, or literature claim.

\section{Why the manuscript has genuine archival value}

The final overview performs several functions that would otherwise be unusually easy to lose.

\begin{itemize}
  \item It reconstructs the historical and technical status of \logic and distinguishes the abstract composition-table semantics from nearby concrete-domain and topological logics whose decidability status is different.
  \item It records the conditional certificate architecture: a machine-checked soundness pipeline around a finite-certificate scheme, together with an explicit statement of the remaining completeness obligation associated with F6.
  \item It surveys the ordered-disjoint RCC5 normal form and the proposed decidability result for the \(\forall\mathsf{PO}\)-free fragment, while marking their precise verification status rather than presenting all claims as equally established.
  \item It preserves failed routes, counterexamples, and repair history: naive blocking, finite quotients, automata, mosaic and certificate variants, width arguments, identity-selector obstructions, and the regular-cover pivot.
  \item It points to prototype reasoners, cross-validation suites, Lean files, finite verification probes, and the result-to-artifact trail needed for reproducibility and future audit.
  \item It states plainly that full decidability remains open, that F6 is unresolved, and that empirical agreement is evidence rather than a completeness proof.
\end{itemize}

This combination is scientifically valuable even without a final resolution. Negative knowledge is especially perishable. Without a citable overview, future researchers are likely to rediscover the same blocking failures, global-coherence gap, width obstruction, and selector examples independently.

\section{What an arXiv posting would and would not claim}

A suitable arXiv posting would make the following claim:

\begin{quote}
\itshape
Here is the present state of a long-standing problem, the verified and unverified components of a sustained investigation, the strongest surviving partial results, the failed routes and their counterexamples, and a reproducible collection of artifacts intended to support future work.
\end{quote}

It would \textbf{not} claim that:

\begin{itemize}
  \item the open satisfiability problem for full \logic has been resolved;
  \item every theorem-level argument has received independent expert review;
  \item the prototype reasoners constitute a proved complete decision procedure for the full logic; or
  \item the empirical test record can replace a completeness theorem.
\end{itemize}

The final manuscript's explicit status labels, open-premise statements, and methodological caveats make this distinction visible enough for responsible dissemination.

\section{Why further delay would be counterproductive}

The overview should not be held until F6 is settled. That could postpone dissemination for years and would defeat the paper's archival purpose. The main value of the document is precisely that it freezes a difficult research landscape at a point where the successful components, failed approaches, formal artifacts, and remaining obligation have been organized into a coherent handoff.

arXiv is well suited to this role. Version~1 can preserve the present state; later versions can incorporate expert review, corrected attributions, stronger formalization, or a future resolution. Posting now does not preclude narrower peer-reviewed papers on the unconditional fragment, the normal form, the formal certificate soundness result, or the computational methodology.

\section{Recommended release package}

I recommend releasing the manuscript after a final mechanical submission check, together with the following archival package:

\begin{enumerate}
  \item the overview PDF and its complete TeX source;
  \item an immutable repository tag or release matching the arXiv version;
  \item a compact \texttt{MANIFEST} or artifact table recording the exact commit, Lean version, entry-point files, verification commands, and status of each principal claim;
  \item a DOI-bearing or otherwise durable archive of the repository snapshot, rather than reliance on a mutable development branch alone; and
  \item citations in the paper to the immutable release or commit, not merely to \texttt{master}.
\end{enumerate}

These measures are important because the overview preserves the intellectual map, while the repository preserves the executable evidence. A future reader needs both.

\section{Two optional presentational improvements}

Neither point below is a condition for arXiv submission.

\textbf{A more technically searchable title.} A title that leads with the logical problem may improve discovery by description-logic and qualitative-spatial-reasoning researchers. One possible form is:

\begin{quote}
\emph{Satisfiability in \logic after Twenty-Five Years: Certified Components, a Decidable Fragment, and the F6 Width Barrier.}
\end{quote}

The agentic-AI question could remain as a subtitle. The present title is engaging, but it may cause indexing and readers to classify the paper primarily as an AI-methodology essay rather than as the main technical handoff for \logic.

\textbf{The DL 2026 rejection appendix.} The appendix is legitimate contextual material, but it is not needed to understand the logical problem, the work packages, or F6. Moving it to a separately linked note or open letter could keep the main overview more focused and reduce the risk that reception of the technical contribution becomes entangled with the dispute. This is a strategic suggestion, not a condition of scientific worthiness.

\section{Assessment against a reasonable arXiv threshold}

With the final corrections in place, the manuscript satisfies a reasonable threshold for an arXiv research report:

\begin{itemize}
  \item it addresses a significant, long-open problem;
  \item it provides a structured state-of-the-art synthesis;
  \item it reports nontrivial certified, theorem-level, and empirical partial results with differentiated status labels;
  \item it isolates unresolved proof obligations rather than concealing them;
  \item it preserves reusable implementations and formal artifacts;
  \item it documents failed approaches that materially narrow future work; and
  \item it offers a transparent account of attribution, AI assistance, and the limits of the evidence.
\end{itemize}

The paper's value is therefore not merely that publication is preferable to losing the files. It is a coherent scientific record whose durable availability can reduce duplicated effort and improve the starting position of future researchers.

\section{Final recommendation}

\fboxtext{%
\textbf{Post the overview as arXiv version 1 after the final mechanical and artifact-link checks.} Preserve a matching immutable repository snapshot. Present the work explicitly as a research overview and technical handoff, not as a solution of the open problem. Peer-reviewed extraction of individual results can follow later and should not be a precondition for archiving the overview.}

My confidence in this dissemination recommendation is high. The recommendation does not independently validate every mathematical claim or the manuscript's July 2026 literature survey; it judges whether the carefully qualified overview, as supplied, is a scientifically useful and responsible object for public archival dissemination.

\vfill
\begin{flushright}
\textit{GPT-5.6 Pro (OpenAI)}\\
20 July 2026
\end{flushright}

\end{document}
